\section{A Bernoulli unit excludes the ordinary rank-two source}
\label{sec:kt-ordinary-rank-two-unit}

This section proves the necessary ordinary direction
that was absent from the earlier field-variation
discussion.  It uses the strict Selmer vanishing in
KT-AUXILIARY-UNIT-SELMER-BRIDGE and a stable lattice,
not a converse to an Eisenstein congruence theorem.

\begin{theorem}[KT-ORDINARY-RANK2-UNIT-EXCLUSION]
Let $p>3$, let $D<0$ be a fundamental discriminant
prime to $p$, put $E=\mathbf Q(\sqrt D)$ and
$\eta=\chi_D$, and let $a>0$ be even with
$a\not\equiv0\pmod{p-1}$.  Fix a finite coefficient
field $K/\mathbf Q_p$, with valuation ring $O$ and
residue field $k$.  There is no continuous
absolutely irreducible representation
\[
                 r:G_{\mathbf Q}\longrightarrow
                                  \mathrm{GL}_2(K)        \tag{OR.1}
\]
satisfying all the following conditions if
$B_{1,\eta\omega^a}$ is a $p$-adic unit:
\begin{enumerate}
\item Its residual semisimplification is
      $1\oplus\bar\eta\bar\varepsilon_p^a$.
\item Its restriction to $G_{\mathbf Q_p}$ has a
      stable line whose rank-one quotient is
      unramified.
\item Its restriction to $G_E$ is unramified at
      every finite place away from $p$.
\end{enumerate}
The conclusion holds whether $p$ splits or is inert
in $E$, and it places no upper bound on $a$.
\end{theorem}

\begin{proof}
Write $\delta=\bar\eta\bar\varepsilon_p^a$.
The residual characters $1,\delta$ are distinct
already on inertia at $p$, since $p\nmid D$ and
$a\not\equiv0\pmod{p-1}$.
Apply Urban's stable-lattice theorem in the order
$\rho_1=1$, $\rho_2=\eta\varepsilon_p^a$, with ideal
the coefficient maximal ideal.  More precisely use
Theorem~1.1 of \emph{Selmer Groups and the
Eisenstein--Klingen Ideal}, arXiv
\texttt{math/9803171v1}, author p.~6.
The image of the continuous representation is
compact and admits a stable lattice, so its
characteristic polynomials are integral; their
residual congruence is hypothesis~1.  The two
residual blocks are absolutely irreducible and
distinct, and (OR.1) is absolutely irreducible.
All hypotheses of that theorem hold.

It gives a stable lattice $L$ with a nonsplit
residual extension
\[
 0\longrightarrow k\longrightarrow\bar L
       \longrightarrow k(\delta)\longrightarrow0.        \tag{OR.2}
\]
Its auxiliary fractional module is free of rank
one over the valuation ring, so the left term is
exactly one copy of $k$.  The theorem's condition
that $L$ has no quotient isomorphic to the first
residual block makes (OR.2) nonsplit.  Its class is
therefore nonzero in
\[
             H^1(\mathbf Q,k(\bar\eta\omega^{-a})).       \tag{OR.3}
\]

Let $V^+\subset K^2$ be the ordinary stable line
in hypothesis~2.  Its intersection with $L$ is
saturated, and the quotient lattice is torsion-free.
Hence reduction gives a stable line and its
rank-one quotient in $\bar L|_{G_{\mathbf Q_p}}$.
The quotient is unramified.  Of the two residual
characters $1,\delta$, only $1$ is unramified on
inertia at $p$, so this quotient has character $1$
and the saturated line has character $\delta$.
That line is distinct from the $k$-line in (OR.2).
Its map to the quotient $k(\delta)$ is consequently
an isomorphism, splitting (OR.2) on the full local
decomposition group.  Thus (OR.3) is strict at $p$.

For every $q\ne p$, hypothesis~3 makes
$I_q\cap G_E$ act trivially on $r$, hence also
on (OR.2).  The latter representation on $I_q$
factors through
\[
                     I_q/(I_q\cap G_E),                 \tag{OR.4}
\]
which has order at most two.  Since $p>2$,
averaging splits every extension over this finite
quotient.  Therefore the restriction of (OR.3)
to $H^1(I_q,k(\bar\eta\omega^{-a}))$ is zero.
This argument includes the ramified primes $q\mid D$:
it does not incorrectly treat $\eta|_{I_q}$ as
trivial there.

The class (OR.3) is now in precisely the residual
strict group of KT-AUXILIARY-UNIT-SELMER-BRIDGE,
with vanishing on $G_{\mathbf Q_p}$ and on every
away-prime inertia group.  That theorem makes
the group with coefficients
$\mathbf F_p(\bar\eta\omega^{-a})$ zero under the
Bernoulli-unit hypothesis.  Extending scalars to
$k$ preserves this statement: these character
values lie in $\mathbf F_p$, and a choice of
$\mathbf F_p$-basis of $k$ identifies the coefficient
module and all local maps with a finite direct
sum of the original ones.  Thus (OR.3) must be
zero, contradicting its construction.
\end{proof}

\begin{corollary}[KT-ORDINARY-RANK2-RESTRICTION-VARIANT]
The first hypothesis in the theorem may be replaced
by the two conditions
\[
 \bigl(\bar r|_{G_E}\bigr)^{\rm ss}
                  =1\oplus\bar\varepsilon_p^a,
 \qquad \det\bar r=\bar\eta\bar\varepsilon_p^a.          \tag{OR.5}
\]
\end{corollary}
\begin{proof}
The residual image is finite.  For each simple
module for that finite image, its socle on the
normal subgroup coming from $G_E$ is nonzero and
stable under the full group, hence is the whole
module.  Thus restriction of the semisimplified
residual representation to $G_E$ is semisimple.
The two characters on $G_E$ are distinct because
$p\nmid D$ makes $E$ disjoint from
$\mathbf Q(\mu_p)$, and the exponent is nonzero
modulo $p-1$.  They are invariant under conjugation
by $G_{\mathbf Q}$.  In the semisimplified residual
representation their individual eigenspaces are
therefore stable under $G_{\mathbf Q}$.  The two
characters over $\mathbf Q$ must be
$\bar\eta^b$ and
$\bar\eta^{1-b}\bar\varepsilon_p^a$ for some
$b\in\{0,1\}$, by the determinant condition.
Replace $r$ by $r\otimes\eta^b$.
It remains absolutely irreducible; its restriction
to $G_E$ is unchanged, and its ordinary quotient
remains unramified at $p$ because $p\nmid D$.
Its residual pair now has the form in the theorem,
which gives the contradiction.
\end{proof}

\paragraph{Arithmetic use and remaining scope.}
For the original even parameter $e$ with GD,
both $a=e$ and $a=2e$ satisfy the nontrivial
inertia hypothesis.  The corresponding auxiliary
units therefore exclude the ordinary rank-two
sources with these residual pairs and with the
stated unramified restriction to $E$.
Equivalently, any such source forces the relevant
$B_{1,\eta\omega^a}$ to be a nonunit.
The integrality used in that formulation is also
contained in the Stickelberger calculation of
the auxiliary-unit bridge, whose auxiliary
multiplier is a $p$-adic unit.

When the source is described by classical forms
of weight $a+1$ and quadratic nebentypus $\eta$,
the determinant in (OR.5) is the usual residual
determinant $\bar\eta\omega^a$.  The ordinary and
unramified-after-$E$ hypotheses must still be
matched to the actual source's local predicates;
they are explicit hypotheses of the result above.
No ordinary condition is silently imposed on
an arbitrary cuspidal form.  No generic stable
rank-three source is excluded by this theorem.

\begin{corollary}[KT-AUXILIARY-UNIT-NO-TYPEII]
Retain the carried source convention and take actual
source data $\mathcal D$ for the imaginary quadratic
field $E$, with $N=2$, $\psi=1$, even $e\ge2$, and
GD.  Let $S$ be its normalized global source ring,
let $f:S.R\to\overline{\mathbf Q}_p$ be an actual
\texttt{IsLocPt} specialization, and let
$r=\operatorname{push}_f(S.r)$.
If
\[
                     B_{1,\eta\omega^{2e}}
                              \in\mathbf Z_p^\times,     \tag{OR.6}
\]
then $r$ is not of Type~II.
This is an exclusion of actual source points;
the bare \texttt{Adm2A} predicate alone is not
being asserted to include residual Eisenstein
data or the determinant $\eta$.
\end{corollary}

\begin{proof}
First the point has finite coefficients.
The literal private theorem \texttt{fin\_quot}
in \texttt{UnitaryThreeGroups/TanDet.lean:212--216}
proves that $S.R/(\ker f)^2$ is finite over the
base coefficient valuation ring.  Its quotient
$S.R/\ker f$ is therefore finite as well.  This
quotient embeds in $\overline{\mathbf Q}_p$, so
its fraction field is a finite extension of the
base coefficient field and contains every matrix
entry of $r$.  The hypotheses are exactly
\texttt{IsCO} and \texttt{IsLocPt}; no
torsion-freeness of the entire source ring or
automorphic descent is used.

Suppose $r$ has Type~II.  The carried
\texttt{typeII\_facts} theorem supplies a
continuous absolutely irreducible representation
$W$ with $\det W=\eta$ and matrices
\[
 W(g)=
 \begin{cases}
  \operatorname{corner}(r(g).\mathrm{mat}),&\eta(g)=1,\\
  \operatorname{corner}(r(g).\mathrm{mat})
                   \operatorname{diag}(-1,1),&\eta(g)=-1.
 \end{cases}                                               \tag{OR.7}
\]
These matrices belong to the same finite
coefficient field.  Enlarge it by the finitely
many entries of a local flag matrix when needed.
The exact source locators are
\texttt{CoboundaryLemmaNoColumn.lean:135--141}
for the properties of $W$ and
\texttt{BaseRepGL3NormalisedOfPsiEqOne.lean:28--34}
for its matrix formula.

The entrywise residual assertion
\texttt{IsPtOf.isPt\_push:124--145}, or the
\texttt{red} field in
\texttt{RhoKApply.lean:76--83}, says that
$r(g).\mathrm{mat}$ reduces to
$\operatorname{diag}(\bar\chi^e,1,\bar\chi^{-e})$.
Formula (OR.7) therefore gives the actual global
residual matrix
\[
       \bar W(g)=
           \operatorname{diag}(\bar\eta\bar\chi^e,
                                      \bar\chi^{-e}).     \tag{OR.8}
\]
Consequently $R=W\otimes\chi^e$ has residual
semisimplification $\bar\eta\bar\chi^{2e}\oplus1$.
For this canonical $W$, no extra quadratic twist
is required.

The whole local argument in
\texttt{AdmissibleOfTypeII.lean:142--195}
constructs a flag on the full decomposition group
with ordered inertial characters
$(\chi^e,\chi^{-e})$, since $\psi=1$.
Thus $R$ has ordered inertia $(\chi^{2e},1)$
and an unramified quotient.  This statement
also holds at an inert prime; the proof uses
the full rational decomposition group.

Finally \texttt{PtInertiaAway.lean:54--69}
shows that on $I_q\cap G_E$ the source point is
conjugate to the base inertia representation.
For $q\ne p$ and $\psi=1$ the latter is the
identity.  Hence $W$, and also $R$, is
unramified away from $p$ after restriction to
$E$.  GD gives $(p-1)\nmid2e$ and $p\nmid D$.
All hypotheses of
KT-ORDINARY-RANK2-UNIT-EXCLUSION now hold for
$a=2e$, contradicting (OR.6).
\end{proof}

The actual GD hypotheses are retained here.  In
particular with $N=2$ they require odd auxiliary
discriminant, as explicitly derived in
\texttt{AdmissibleOfTypeII.lean:119}.
The generic rank-two theorem itself has no such
odd-discriminant restriction.  The unit for $a=e$
likewise excludes any ordinary lower-weight
rank-two seed satisfying its stated residual
and unramified-after-$E$ conditions; this includes
the construction underlying the ordinary
symmetric-square examples when those hypotheses
hold.  Neither assertion eliminates an arbitrary
irreducible stable rank-three source.
